\section{讨论}

\subsection{架构优势对比分析}

本研究通过对比纯CNN、纯ViT、串行融合（Hybrid）、并行融合（Parallel）与嵌入式注意力（Embedded）五种架构，总结了各自的优势与局限。纯CNN依托局部感受野与权重共享，在小样本场景下收敛稳定、参数效率高，但难以捕捉跨区域的占位效应。纯ViT在理论上具备灵活的全局建模能力，然而缺少卷积带来的归纳偏置，导致在本数据集上整体性能与召回率均显著低于其他模型。

\subsubsection{Pure CNN模型的优势与局限}

CNN对纹理与边界的敏感性使其在脑膜瘤与无肿瘤类别上获得高召回，同时保持较低的计算开销；局部感受野限制了其对胶质瘤浸润性变化的捕获，召回率仅24.00\%，提示需要额外的全局信息补充。

\subsubsection{Pure ViT模型的优势与局限}

ViT依赖自注意力在单层内建立远程依赖，能够理论上覆盖跨区域语义关系；但在缺乏大规模预训练的情况下，训练曲线波动较大，四个类别的F1均落后于卷积基线，说明对数据规模与正则化的依赖更高。

\subsubsection{串行融合（Hybrid）}

Hybrid通过“CNN特征提取+Transformer编码”串行组合，在不显著增加计算量的前提下提升全局感知能力，整体准确率达到76.40\%，在胶质瘤等难例的召回上略优于纯CNN，但仍受限于串行路径对全局信息的表达容量。

\subsubsection{并行融合（Parallel）}

Parallel在并行支路中同时利用卷积的局部敏感性与Transformer的全局注意力，在特征层融合后输出决策，是本次实验的最佳方案（准确率80.71\%、F1 77.85）。并行结构显著提升胶质瘤召回至35.00\%，并降低无肿瘤的误报，验证了局部—全局互补的有效性。

\subsubsection{嵌入式注意力（Embedded）}

Embedded在卷积骨干中插入注意力瓶颈，实现轻量的上下文注入。该设计几乎不改变整体流程，却在脑膜瘤、无肿瘤和垂体瘤类别上与并行模型相当，为资源受限场景提供了兼顾精度与复杂度的折中方案。

\subsection{与放射科医生诊断流程的对应关系}

融合模型的处理顺序与临床诊断高度一致：医生先关注病灶边界和纹理等局部细节，再评估侧脑室压迫、中线偏移等全局占位效应。并行与嵌入式设计在模型层面同时保留这两类线索，有助于模拟这一推理流程并减少胶质瘤漏检。

\subsection{未来改进方向}

后续工作可在大规模医学影像预训练、注意力可视化与不确定性量化上深化；多模态序列（T1/T2/FLAIR）融合和模型轻量化也是提升临床可用性的关键方向。

\subsection{临床应用前景}

最佳模型能够作为辅助诊断工具或快速初筛系统，为放射科医生提供第二意见，也可用于教学演示典型影像模式；但最终决策仍需专家审阅，模型输出应被视作风险提示而非唯一依据。
